%%%%%%%%%%%%%%%%%%%%%%%%%%%%%%%%%%%%%%%%%%%%%%%%%%%%%%%%%%%%%%%%%%%
%                                                                 %
%                 Packages / Grundeinstellungen                   %
%                                                                 %
%%%%%%%%%%%%%%%%%%%%%%%%%%%%%%%%%%%%%%%%%%%%%%%%%%%%%%%%%%%%%%%%%%%


% Festlegung des Allgemeinen Dokumentenformat
\documentclass[a4paper,12pt,parskip=half,headsepline,DIV=12,numbers=noenddot]{scrartcl}

%%%%%% Muss in die documentclass %%%%%%%%
%BCOR=12mm, Korrektur fuer die Bindung
%DIV12 DIV-Wert fuer die Erstellung des Satzspiegels

% Keine floats in andere Sections
\usepackage[section]{placeins}

% Weitere Pakete
\usepackage{microtype}
\usepackage{caption}
\usepackage{fontspec}
\usepackage{pdflscape}
\usepackage{float}
\usepackage{dirtree}
\usepackage{subcaption}
\usepackage{enumitem}

% Booktabs Tabellen
\usepackage{tabularray}
\UseTblrLibrary{booktabs}
\DefTblrTemplate{contfoot-text}{normal}{Fortsetzung auf nächster Seite}
\SetTblrTemplate{contfoot-text}{normal}
\DefTblrTemplate{conthead-text}{normal}{}
\SetTblrTemplate{conthead-text}{normal}

% Um Captions in Tabellen zu deaktivieren 
%\DefTblrTemplate{caption-tag}{default}{}
%\DefTblrTemplate{caption-sep}{default}{}
%\DefTblrTemplate{caption-text}{default}{}

% Grafiken aus PNG Dateien einbinden
\usepackage{graphicx}

% Deutsche Sonderzeichen und Silbentrennung nutzen
\usepackage[ngerman]{babel}
\usepackage{blindtext}

% Eurozeichen einbinden
\usepackage[right]{eurosym}

% Kopf- und Fußzeilen
\usepackage[headsepline,autooneside=false]{scrlayer-scrpage}
\clearpairofpagestyles

% Schriftart 
\usepackage{lmodern}
\setmainfont{TeX Gyre Termes}
\setsansfont{TeX Gyre Adventor}

% Floatende Bilder ermöglichen
\usepackage{floatflt}

% tikz
\usepackage{tikz}
\usetikzlibrary{calc,arrows,math}
\usetikzlibrary{shapes.geometric,positioning}

%% Schaltpläne nach europäischen Richtlinien
\usepackage[european]{circuitikz}
\tikzset{x=1mm,y=1mm}

\usepackage{siunitx}
\sisetup{output-decimal-marker={,},detect-all}

% Bricht lange URLs "schön" um
\usepackage[hyphens,obeyspaces,spaces]{url}

% Paket für Textfarben
\usepackage{xcolor} 
\definecolor{LightGray}{gray}{0.9}
\usepackage[pagecolor=white]{pagecolor}

% Mathematische Symbole importieren
\usepackage{amssymb}

% Paket für Zeilenabstand
\usepackage{setspace}

% Für Bildbezeichner
\usepackage{capt-of}

% Für Stichwortverzeichnis
\usepackage{makeidx}

% Für if und while 
\usepackage{etoolbox}

% Konfiguriere das Inhaltsverzeichnis
\usepackage{tocbasic}
\DeclareTOCStyleEntries[
  raggedentrytext,
  %numwidth=0pt, if numbers=noenddot is not set
  numsep=1ex,
  dynnumwidth,
]{tocline}{chapter,section}
\DeclareTOCStyleEntries[
  linefill=\TOCLineLeaderFill,
]{tocline}{section,subsection,subsubsection,paragraph,subparagraph}

\newcommand*\tocentryformat[1]{{\rmfamily#1}}
\RedeclareSectionCommands
  [tocentryformat=\tocentryformat,tocpagenumberformat=\tocentryformat]{subsection,subsubsection,paragraph,subparagraph}

\newcommand*\tocentrysectionformat[1]{{\rmfamily\bfseries#1}}
\RedeclareSectionCommands
  [tocentryformat=\tocentrysectionformat,tocpagenumberformat=\tocentrysectionformat]{section}  
  
\DeclareTOCStyleEntries[
  pagenumberbox=\hbox,
  dynnumwidth]{tocline}{chapter,section,subsection,subsubsection,paragraph,subparagraph,figure,table}

% Für schönere Listings
\usepackage[outputdir=log, newfloat,]{minted}
\setminted{
  frame=lines,
  framesep=2mm,
  baselinestretch=1.2,
  bgcolor=LightGray,
  fontsize=\footnotesize,
  linenos,
  breaklines=true,
  breakanywhere=true,
  autogobble,
  tabsize=2
}
\setmintedinline{}

% Keine Floats bei Listings
\newenvironment{code}[2]
  {\captionsetup{type=listing}
  \providecommand{\captiontitle}{#1}
  \providecommand{\labeltitle}{#2}
  \vspace*{0.3cm}
  }
  {
  \vspace*{-0.8cm}
  \caption{\captiontitle}
  \label{\labeltitle}
  \vspace*{0.35cm}
  }
\SetupFloatingEnvironment{listing}{}

% Nummerierung inkl. Section
\usepackage{chngcntr}
\counterwithin{table}{section}
\counterwithin{figure}{section}
\counterwithin{listing}{section}

% Abkürzungsverzeichnis
\usepackage[printonlyused, smaller, withpage]{acronym}

% Erzeugt Inhaltsverzeichnis mit Querverweisen zu den Abschnitten (PDF Version)
\usepackage[bookmarksnumbered,hyperfootnotes=false,hypertexnames=false]{hyperref}
\hypersetup{
    colorlinks=true,
    linkcolor=black,
    filecolor=blue,
    citecolor = black,      
    urlcolor=blue,
}

% Darf erst hier eingebunden werden! 
\usepackage{subfiles}
\usepackage{csquotes}

% Indexerstellung
\makeindex


%%%%%%%%%%%%%%%%%%%%%%%%%%%%%%%%%%%%%%%%%%%%%%%%%%%%%%%%%%%%%%%%%%%
%                                                                 %                    
%                   Definition Zitierstil                         %
%                                                                 %
%%%%%%%%%%%%%%%%%%%%%%%%%%%%%%%%%%%%%%%%%%%%%%%%%%%%%%%%%%%%%%%%%%%

% Zitierung nach IEEE

\usepackage[
backend=biber,
style=ieee,
autocite=inline,
]{biblatex}
\addbibresource{bibtex/hauptdatei.bib}

% Zitierung nach APA

%\usepackage[
%backend=biber,
%style=apa,
%autocite=inline,
%]{biblatex}
%\addbibresource{bibtex/hauptdatei.bib}

\setcounter{biburllcpenalty}{7000}
\setcounter{biburlucpenalty}{8000}
%%%%%%%%%%%%%%%%%%%%%%%%%%%%%%%%%%%%%%%%%%%%%%%%%%%%%%%%%%%%%%%%%%%
%                                                                 %                    
%                    Definition Deckblatt                         %
%                                                                 %
%%%%%%%%%%%%%%%%%%%%%%%%%%%%%%%%%%%%%%%%%%%%%%%%%%%%%%%%%%%%%%%%%%%

% true für Bachelorarbeit / false für Hausarbeit
\newbool{bachelorarbeit}
\setbool{bachelorarbeit}{false}

% Setze Fachbereich
\newcommand{\department}{Fachbereich II \\ Management und Informationssysteme}

% Setze Studiengang
\newcommand{\studyprogram}{Wirtschaftsinformatik B.Sc.}

% Setze Modulname (bachelorarbeit muss false sein)
\newcommand{\modulname}{Vernetzte Systeme}

% Setze Dozent:in (bachelorarbeit muss false sein)
\newcommand{\auditor}{\textbf{Dozent:} \> Prof. Dr. Oliver Radfelder}

% Setze Gutachter:innen (bachelorarbeit muss true sein)
\newcommand{\firstauditor}{\textbf{Erstgutachter:} \> Prof. Dr. Oliver Radfelder}
\newcommand{\secondauditor}{\textbf{Dozentin:} \> Prof. Dr. Karin Vosseberg}

% Setze Titel und Untertitel der Abreit 
\newcommand{\thetitle}{Analyse verschiedener Dockerkonstellationen}

% Setze Autor:in und MatNr.
\newcommand{\theauthor}{Junior Ekane}
\newcommand{\matriculationnumber}{40128}

% Abstand zwischen Name und MatNr. (siehe Deckblatt)
\newcommand{\myspace}{1.0cm}

% Muss in src/basic_structure/deckblatt.tex einkommentiert werden! 

\newcommand{\secondauthor}{\> Florian Quaas \> MatNr. 39952\\}
\newcommand{\thirdauthor}{\> Herman Tsago \> MatNr. 00000\\}
\newcommand{\fourthauthor}{\> Steve Aguiwo II \> MatNr. 40088\\}

% PDF Metadaten
%\hypersetup{pdfinfo={
%Title={\thetitle},
%Author={\theauthor}
%}}

\hypersetup{pdfinfo={
Title={\thetitle},
Author={\theauthor}
}}

%%%%%%%%%%%%%%%%%%%%%%%%%%%%%%%%%%%%%%%%%%%%%%%%%%%%%%%%%%%%%%%%%%%
%                                                                 %                    
%                     Beginn des Inhalts                          %
%                                                                 %
%%%%%%%%%%%%%%%%%%%%%%%%%%%%%%%%%%%%%%%%%%%%%%%%%%%%%%%%%%%%%%%%%%%

%%%%%%%%%%%%%%%%%%%%%%%%%%%%%%%%%%%%%%%%%%%%%%%%%%%%%%%%%%%%%%%%%%%
%  Special Characters:                                            %
%                                                                 %
%             \& \% \$ \# \_ \{ \}                                %
%             \textasciitilde (~)                                 %
%             \textasciicircum (^)                                %     
%             \textbackslash (\)                                  %                    
%      \glqq Text\grqq{} für Anführungszeichen                    %
%%%%%%%%%%%%%%%%%%%%%%%%%%%%%%%%%%%%%%%%%%%%%%%%%%%%%%%%%%%%%%%%%%%

\begin{document}



% Definition Header Sections sollen in der Kopfzeile stehen; Kopfzeile mit Unterstrich
\automark[subsection]{section}
\KOMAoptions{headsepline=true}
%\ihead{Kopfzeile innen}
%\chead{Kopfzeile außen}
\ohead{\headmark}

% Definition footer \pagemark steht für Seitennummer
%\ifoot{Fußzeile innen}
%\cfoot{Fußzeile Mitte}
\ofoot{\pagemark}

% Hier werden die Trennvorschläge inkludiert
%%%%%%%%%%%%%%%%%%%%%%%%%%%%%%%%%%%%%%%%%%%%%%%%%%%%%%%%%%%%%%%%%%%%%%%%%%%%%%%%%%%%%%%%%%%%
%    Hier müssen alle Wörter eingetragen werden, welche Latex von nicht korrekt trennt    %
%                bzw. bei denen man die genaue Trennung vorgeben möchte.                  %
%%%%%%%%%%%%%%%%%%%%%%%%%%%%%%%%%%%%%%%%%%%%%%%%%%%%%%%%%%%%%%%%%%%%%%%%%%%%%%%%%%%%%%%%%%%

\hyphenation{
	Film-pro-du-zen-ten
	Lux-em-burg
	Soft-ware-bau-steins
	zeit-in-ten-siv
}


% Leere Seite am Anfang
%\thispagestyle{empty} % erzeugt Seite ohne Kopf- / Fusszeile
%\mbox{}
%\newpage

% Titelseite 
%%%%%%%%%%%%%%%%%%%%%%%%%%%%%%%%%
%           Deckblatt           %
%%%%%%%%%%%%%%%%%%%%%%%%%%%%%%%%%
\setmainfont{TeX Gyre Adventor}
\thispagestyle{empty}
\begin{figure}[h!]
	\centering
	\includegraphics[width=0.6\textwidth]{src/abbildungen/logoneu.png}
\end{figure}
\begin{center}
	\large{\textbf{\department}}\\
	\large{\textbf{\studyprogram}}\\
	\vspace{1cm}
	\ifbool{bachelorarbeit}{
		\LARGE{\textbf{Bachelorarbeit}}\\
		\large{zur Erlangung des akademischen Grades \\ Bachelor of Science}\\
	}
	{
		\large{\textbf{Modul\\ \modulname}}\\
	}
	\vspace*{\fill}
	\line(1,0){450}\\
	\doublespacing
	\textbf{\Large{\thetitle}}\\
	\line(1,0){450}\\
\end{center}
\vspace*{\fill}
\onehalfspacing
\small{
\begin{flushleft}
	\begin{tabbing}
		\textbf{Vorgelegt von:} \hspace*{0.8cm}\= \theauthor \hspace*{\myspace}\= MatNr. \matriculationnumber \\


		%%%%%%%%%%%%%%%%%%%%%%%%%%%%%%%%%%%%%%%%%%%%%%
		%                                            %
		%  Hier weitere Autor:innen einkommentieren  %
		%   Müssen hauptdatei.tex definiert sein     %
		%	                                         %
		%%%%%%%%%%%%%%%%%%%%%%%%%%%%%%%%%%%%%%%%%%%%%%		
		\secondauthor
		\thirdauthor
		\fourthauthor
		%\fifthauthor
    \\


		\textbf{Vorgelegt am:} \> \today\\
		\\
    \ifbool{bachelorarbeit}{
			\firstauditor\\
			\secondauditor\\
		}
		{
			\auditor\\
      \secondauditor\\
		}
	\end{tabbing}
\end{flushleft}}
\setmainfont{TeX Gyre Termes}

\newpage

% Singlespacing (Zeilenabstand) (Default)
\singlespacing
\normalsize

% Abstract falls gewünscht
%\thispagestyle{empty}
%\input{abstract}
%\newpage

% Inhaltsverzeichnis anzeigen
\pagestyle{empty}
\tableofcontents
\newpage
\pagestyle{headings}

% Header für den Inhalt 
\KOMAoptions{headsepline=true}
\ohead{\headmark}

% Input Inhalt
\section{Einleitung}\label{einleitung}

\subsection{Vorstellung des Moduls Vernetzte Systeme}

\section{Projektablauf}\label{Projektablauf}

\subsection{Vorstellung der Projektidee}

\subsection{Projektorganisation und Teamarbeit}


Für unsere Auseinandersetzung zum Thema "\/Analyse verschiedener Dockerkonstellationen"\/ im Rahmen der Veranstaltung \/"Vernetzte Systeme"\/ haben wir uns zegielt für einen strukturierten Projektplan, der sich an agilen Methoden orientiert. Zum Entwerfen und Fertigstellung haben wir einen Sprint von drei Wochen gemacht inklusiv das Verfassen vom Bericht. Regelmäßig haben wir uns alle zwei Tage getroffen und da wurde immer über die Fortschritte gesprochen und neue Aufgabe verteilt. Die ersten Tage habe wir erstmal dürber ausgetauscht, was wir machen wollen und wie wir es tun möchten?

\\

Die Implementierung haben wir in mehreren iterativen Phasen unterteilt, wobei das Kanban-System zur Visualisierung und Verwaltung des Arbeitsflusses eingesetzt wurde. Hierfür haben wir die Webanwendung Trello(Verlinkung fehlt)  verwendet. In der App haben wir über einen Projekt-Dashboard verfügt, auf dem unsere Aufgaben in drei Kategorien aufgelistet waren wie:

\begin{itemize}
\item Zu erledigen
\item In Arbeit
\item Erledigt
\end{itemize}

Dies ermöglichte eine flexible Anpassung an wechselnde Anforderungen und eine kontinuierliche Verbesserung der Implementierung. Die agilen Prinzipien unterstützten uns dabei, schnell auf Herausforderungen zu reagieren und eine effektive Implementierung sicherzustellen.

\section{Docker}\label{Docker}


\subsection{Was ist ein Docker?}


\subsection{Was ist ein Netzwerk?}


\subsection{Was ist ein Image?}

Images sind schreibgeschützte Templates mit Anweisungen zum Erstellen eines Containers. Ein docker Image erstellt Container zur Ausführung auf der Docker-Plattform. ein Docker-Image führt Code in einem Docker-Container aus und es ist möglich mehrere Container aus demselben Image zu erstellen.

\subsubsection{Wie erstellt man ein Docker-Image ?}

Um ein Docker-Image zu erstellen, verwendet man den \texttt{docker image build} Befehl zusammen mit weiteren Optionen zur Spezifikation des Build-Kontexts und des Tags des Images. Zum Beispiel:

\begin{verbatim}
docker image build --progress=plain -t image-fancy context
\end{verbatim}

Dieser Befehl erstellt ein Image mit dem Namen \texttt{image-fancy} aus dem Build-Kontext \texttt{context}, wobei \texttt{--progress=plain} eine einfachere Ausgabe während des Build-Prozesses erzeugt.

\subsection{Was ist ein Container ? }

Ein Container ist eine Abstraktion aud fer App-Ebene, die Code und Abhängigkeiten zusammenfasst. Mehrere Container könnenauf demselben Computer ausgeführt werden und den Betriebssystemkernel mit anderen Containern teilen, die jeweils als isolierte Prozesse im Benutzerbereich ausgeführt weden. Container nehmen weniger Platz ein als VMs und können mehr Anwendungen verarbeiten.

\subsubsection{Wie erstellt und startet man einen Container?}

Ein Container wird basierend auf einem Docker-Image erstellt und gestartet. Das folgende Beispiel zeigt, wie ein Container mit dem Namen \texttt{services} erstellt und konfiguriert wird:

\begin{verbatim}
docker container create \
    --name "services" \
    --hostname "services" \
    --network mynet \
    --volume "$PWD/www/:/var/www/html" \
    --volume "$PWD/log/:/log" \
    --publish 80:80 \
    --publish 8080:8080 \
    "image-fancy"
\end{verbatim}

Dieses Kommando erstellt einen Container, der zum \texttt{mynet} Netzwerk gehört, Volumes für Webinhalte und Logdateien einbindet und Ports 80 sowie 8080 veröffentlicht.

Um einen Docker-Container zu starten, verwendet man den Befehl \texttt{docker container start} gefolgt vom Namen oder der ID des Containers, den man starten möchte. Dieser Befehl ändert den Zustand des Containers von gestoppt zu laufend.

Wir haben einen Docker-Container mit dem Namen \texttt{services} erstellt. Um diesen Container zu starten, der Befehl dazu lautet:

\begin{verbatim}
docker container start services
\end{verbatim}

Dieser Befehl sucht nach einem Docker-Container, der \texttt{services} heißt, und startet ihn. Sobald der Container gestartet ist, läuft er im Hintergrund weiter, und die in ihm enthaltene Anwendung beginnt mit ihrer Ausführung basierend auf den beim Erstellen des Containers festgelegten Konfigurationen und Einstellungen.

\subsubsection{Überprüfung des Container-Status}

Um zu überprüfen, ob der Container erfolgreich gestartet wurde und läuft, kann man den Befehl \texttt{docker container ls} verwenden, der alle laufenden Container auflistet. Wenn Sie auch gestoppte Container sehen möchten, können Sie \texttt{docker container ls --all} hinzufügen, um eine vollständige Liste zu erhalten.

\subsection{Wie stoppt und entfernt man einen Container?}

Zum Stoppen und Entfernen eines Containers verwendet man die Befehle \texttt{docker container stop} und \texttt{docker container rm}. Das folgende Skript demonstriert diesen Prozess für einen Container namens \texttt{services}:

\begin{verbatim}
if docker container ls | grep -q "\bservices$"; then
    docker container stop services
fi

if docker container ls --all | grep -q "\bservices$"; then
    docker container rm services
fi
\end{verbatim}

Zusätzlich entfernt das Skript veraltete SSH-Konfigurationen für den Container aus \texttt{\~/.ssh/docker\_config}.

\subsection{Was sind die Vor - und Nachteile von Docker ? }

\subsubsection{Vorteile}

Docker bietet eine Vielzahl von Vorteilen, die es zu einem beliebten Werkzeug in der Softwareentwicklung und im IT-Betrieb gemacht haben:

\begin{itemize}
    \item \textbf{Portabilität:} Einmal erstellte Docker-Images können über verschiedene Umgebungen hinweg ohne Änderungen ausgeführt werden. Dies erleichtert die Bereitstellung und Migration von Anwendungen erheblich.
    \item \textbf{Konsistenz:} Docker-Container bieten eine konsistente Umgebung für die Anwendung, unabhängig davon, wo der Container ausgeführt wird. Dies reduziert das Problem "es funktioniert auf meinem Rechner nicht", indem Diskrepanzen zwischen Entwicklung, Test und Produktion minimiert werden.
    \item \textbf{Isolation:} Container sind voneinander isoliert und verpacken ihre eigene Software, Bibliotheken und Konfigurationen. Dies verbessert die Sicherheit und vereinfacht Abhängigkeitsmanagement.
    \item \textbf{Ressourceneffizienz:} Docker nutzt die Ressourcen des Host-Betriebssystems effizienter als traditionelle Virtualisierungstechnologien, was zu einer besseren Hardwareauslastung führt.
    \item \textbf{Schnelle Bereitstellung:} Container können innerhalb von Sekunden gestartet und gestoppt werden, was die Bereitstellungszeiten verkürzt und eine schnelle Skalierung ermöglicht.
    \item \textbf{Entwicklungseffizienz:} Docker unterstützt Entwickler mit der Möglichkeit, lokale Umgebungen schnell einzurichten, was die Entwicklungszyklen beschleunigt.
    \item \textbf{Mikroservice-Architekturen:} Docker eignet sich hervorragend für die Entwicklung und den Betrieb von mikroservice-basierten Anwendungen, indem es die Unabhängigkeit und Isolation der Services fördert.
\end{itemize}

\subsubsection{Nachteile}

Trotz seiner vielen Vorteile hat Docker auch einige Nachteile, die berücksichtigt werden sollten:

\begin{itemize}
    \item \textbf{Sicherheitsrisiken:} Die Isolation zwischen Containern ist nicht so stark wie bei vollständig virtualisierten Umgebungen, was potenzielle Sicherheitsrisiken birgt. Es ist wichtig, Sicherheitspraktiken wie das Prinzip der geringsten Privilegien und regelmäßige Updates zu befolgen.
    \item \textbf{Komplexität:} Die Verwaltung von Docker-Containern und -Orchestrierung, insbesondere in großem Maßstab, kann komplex sein und erfordert ein gutes Verständnis von Docker-Netzwerken, Volumes und Sicherheitskonfigurationen.
    \item \textbf{Speicher- und Datenmanagement:} Persistente Daten und deren Verwaltung können in containerisierten Anwendungen eine Herausforderung darstellen, besonders wenn es um Zustandsbehaftete Dienste wie Datenbanken geht.
    \item \textbf{Plattformabhängigkeiten:} Obwohl Docker auf Portabilität ausgelegt ist, können bestimmte Images aufgrund von Systemabhängigkeiten (z.B. Kernel-Versionen) nicht auf allen Plattformen ausgeführt werden.
    \item \textbf{Performance-Überlegungen:} Während Container im Allgemeinen effizient sind, können spezifische Workloads oder Konfigurationen in Containern im Vergleich zu nativen Ausführungen oder anderen Virtualisierungstechnologien Leistungseinbußen erleiden.
    \item \textbf{Lernkurve:} Die umfangreichen Funktionen und das Ökosystem von Docker können für Neueinsteiger überwältigend sein und erfordern Zeit und Ressourcen, um effektiv genutzt zu werden.
\end{itemize}

\subsection{Wie haben wir Docker verwendet?}

In unseren Skripten haben wir Docker verwendet, um eine isolierte und konsistente Umgebung für unsere Anwendungen zu schaffen. Die Skripte automatisieren den Prozess des Builds, der Bereitstellung und des Managements von Containern. Hier einige Beispiele:

\begin{itemize}
    \item \textbf{Image-Erstellung:} Wir nutzen den \texttt{docker image build} Befehl, um Docker-Images aus unseren Anwendungsquellen zu erstellen. Dies ermöglicht eine einheitliche und reproduzierbare Umgebung für die Entwicklung und das Deployment unserer Anwendung.
    \item \textbf{Container-Management:} Unsere Skripte automatisieren das Erstellen, Starten, Stoppen und Entfernen von Containern. Dies vereinfacht die Verwaltung der Lebenszyklen unserer Anwendungen erheblich.
    \item \textbf{Netzwerk und Volumen:} Wir konfigurieren Netzwerke und binden Volumes ein, um die Kommunikation zwischen Containern zu ermöglichen und persistente Daten zu verwalten. Dies ist entscheidend für komplexe Anwendungen, die aus mehreren Diensten bestehen und Zustandsdaten speichern müssen.
\end{itemize}


\section{Ergebnis}\label{Ergebnis}

%Hier später bei Fortschritt irgendwas hinzufügen

\subsection{Etwas}
\subsubsection{noch etwas}
\subsubsection{noch etwas}
\subsection{Etwas}
\subsubsection{noch etwas}
\subsubsection{noch etwas}

\section{Lasttest}\label{Lasttest}


\subsection{Was haben wir analysiert?}



\subsection{Wie haben wir das analysiert?}



\subsection{Endergebnis}

\section{Reflexion}\label{Reflexion}

\section{Fazit}\label{Fazit}



\newpage

% Literaturverzeichnis anzeigen
\ohead{Literaturverzeichnis} % Korrektur für Header 
\phantomsection
\addcontentsline{toc}{section}{Literaturverzeichnis}
\renewcommand\refname{Literaturverzeichnis}
\printbibliography
\newpage



% Abkürzungsverzeichnis anzeigen
%\ohead{Abkürzungsverzeichnis} % Korrektur für Header 
%\section*{Abkürzungsverzeichnis}
%%%%%%%%%%%%%%%%%%%%%%%%%%%%%%%%%%%%%%%%%%%%%%%%%%%%%%%%%%%%%%%%%%%%%%%%%%%%%%%%%%%%%%%%%%
%                             Mit Einrückung im Verzeichnis                             %                     
%  Alle Abkürzungen müssen mit: \acro{acronym}[shortname]{fullname} eingetragen werden. %
%                         Im Fließtext mit \ac{acronym} aufrufen.                       %
%                 Erster Aufruf: fullname (shortname), danach shortname.                %                                   
%%%%%%%%%%%%%%%%%%%%%%%%%%%%%%%%%%%%%%%%%%%%%%%%%%%%%%%%%%%%%%%%%%%%%%%%%%%%%%%%%%%%%%%%%

\begin{addmargin}[1.5em]{0em}

    \begin{acronym}[HTTPS] %längster shortname muss eingetragen werden.

        \acro{eugh}[EuGH]{Europäischer Gerichtshof}

        \acro{eu}[EU]{Europäische Union}

        \acro{https}[HTTPS]{Hypertext Transfer Protocol Secure}

    \end{acronym}

\end{addmargin}
%\addcontentsline{toc}{section}{Abkürzungsverzeichnis}
%\newpage



% Selbstständigkeits Erklärung
\phantomsection
\addcontentsline{toc}{section}{Selbstständigkeitserklärung}

% Header für Erklärung
\ohead{Selbstständigkeitserklärung}

% Input Erklärung
\section*{Selbstständigkeitserklärung}

Wir versichern, die von uns vorgelegte Arbeit selbstständig verfasst zu haben. Alle Stellen, die wörtlich oder sinngemäß aus veröffentlichten oder nicht veröffentlichten Arbeiten anderer entnommen sind, 
haben wir als entnommen kenntlich gemacht. Sämtliche Quellen und Hilfsmittel, die wir für die Arbeit benutzt haben, sind angegeben. Die Arbeit haben wir mit gleichem Inhalt bzw. in wesentlichen 
Teilen noch keiner anderen Prüfungsbehörde vorgelegt.


%\begin{tblr}{X[l] X[l]}
Bremerhaven, den \today & Unterschrift:\\
 % \textit{Junior Ekane}\\
 % \textit{Florian Quaas}\\
  %\textit{Steve Aguiwo}\\
 % \textit{Herman Tsago}
%\end{tblr}


% Leere Abschlussseite
%\newpage
%\thispagestyle{empty} % erzeugt Seite ohne Kopf- / Fusszeile
%\mbox{}

\end{document}
